\labeledsubsubsection{Agent loading}
The plugin-system needs a way of loading a plugin.
In \texttt{Java} loading means either, appending it to the classpath of the current environment (JRE\footnote{Java-Runtime-Environment}), or, creating a new secondary classpath with a new class loader implementation.

The official way of modifying the class-path at runtime is to use an \texttt{Java Agent} / the \texttt{java.lang.Instrument} package. \\
From the \textcite{javadoc:instrument}:
\begin{quote}
    [The package p]rovides services that allow Java programming language agents to instrument programs running on the JVM\footnote{Java-Virtual-Machine}.\newline
    The mechanism for instrumentation is modification of the byte-codes of methods.
\end{quote}~\parencite{javadoc:instrument}\\
This is usually used by diagnostic tools, such as debugger, tracer or profiler, to get access over the current JRE and heap to perform modifications or analytic work.
However, the experiment will use it to inject our dependency (plugin) into the classpath and then access the class loader to find the injected classes and finally invoke them.

For an \texttt{agent} to work properly, certain preparations have to be made:
First, a \lstinline{agentmain} and \lstinline{premain} has to be defined.
Next, the classpath of those methods must go into the \texttt{MANIFEST.MF}.
Finally, the application must be run with a special VM parameter to accept agents.

The following code shows how to implement those methods.
\lstinputlisting[language=Java, linerange={36-41, 264-290}]{./code/plugin-system/src/main/java/nl/fontys/sear/plugin_system/system/PluginLoader.java}
The \lstinline{premain} gets called when the JRE loads the agent JAR.
The \lstinline{agentmain} gets called when the \texttt{agent} got attached to the process.
In most cases one of them is enough, however as \texttt{agent} implementation differs from different Java versions it is safer to define both and let them call each other.
Both methods have two arguments: 
The \lstinline{premain}/\lstinline{agentmain} arguments and the \lstinline{instrumentation}.
The plugin-system needs the \lstinline{instrumentation}, thus it has to be stored globally/statically somewhere, where the plugin-system later can access it.
It is important to \textbf{not} invoke anything else here and \textbf{not} call the \textit{standard} \lstinline{main}.

After defining both methods, the classpath of the class(es) these methods are in must be defined in the \lstinline{MANIFEST.MF}.\\
From the experiment:
\lstinputlisting[
    language=Manifest,
    showspaces=false,
    showtabs=false,
    breaklines=true,
    showstringspaces=false,
    breakatwhitespace=true,
    numbers=left,
    basicstyle=\footnotesize\ttfamily\color{darkgray},
    keywordstyle=\bfseries\color{darkgray},
    commentstyle=\itshape\color{darkgray},
    identifierstyle=\color{darkgray},
    stringstyle=\color{darkgray},
]{./code/plugin-system/src/main/resources/META-INF/MANIFEST.MF}

Finally, when starting the application it is important to add the \lstinline{-javagent:<PATH TO JAR>} as a VM parameter.
For simplicity reasons we put the \lstinline{premain} and \lstinline{agentmain} into the same JAR.
However, it is also possible to put them in a separate JAR.
Either way, the location of this JAR must be set.
Otherwise, the JVM will not recognize the \lstinline{premain}, nor \lstinline{agentmain}, and the \lstinline{instrumentation} variable will be \texttt{null}.